% Pursuing Stacks (cotes 134-2 to 134-8): archive cote 134-7, site page 59 (PDF page 60), typescript page 501.
% typescript. Draft compared: ps6.tex:2907-2961.
% AI-transcribed by Claude Opus 5.5 (claude-opus-5-5, adaptive thinking, effort high), Message Batches API, batch ps4 (msgbatch_012Aoc1NuskATPsyBeKSYK7J), 2026-11-12; not yet checked by eye.
\scanpage{134-7/59}
\mspage{501}
entirely conclusive maybe,  but still it was enough to shake my confidence
that the assumption was OK.  T\unsure{\hand{h}}e counterexample is in terms of cohomology
\[ H^2(1,\underline{Z};\ \underline{Z}) = H^2(\underline{Z},\underline{Z}) = H^2(B_{\underline{Z}},\ \underline{Z}) \]
rather than homology, as usual.   As the classifying space $B_{\underline{Z}} = S^1$ of $\underline{Z}$
is one dimension\hand{al,} it's cohomology is zero in dimensions $i \hand{$\geq$} 2$.  On the \xout{oth}
other hand,  $H^2$ classifies central extensions of $\underline{Z}$ by $\underline{Z}$ ,  and an\deleted{direct}
immediate direct argument shows \deleted{too}\inserted{indeed} that such extensions (indeed, any exten%
\hand{-}sion of $\underline{Z}$ by any group) split.  If we take the schematic $H^2$,  defined by
Hochschild cochains which are polynomial functions,  we get the classifica-%
tion of central extensions of $\underline{G}_a$ \deleted{\xout{\illegible{1}}}by $\underline{G}_a$ \xout{\illegible{1}} ,  as group schemes over
the ring of integers.  Now it is e\hand{a}sy to find such an extension which does
\uline{not} split,  the first one one may think of being the group scheme represen-%
ting the functor
\[ k \hand{$\mapsto$} W_2(k) = \begin{array}[t]{l} \xout{\text{ring}}\inserted[typed]{\text{group}}\text{ of truncated power series }\ 1+at+bt^2 \\ \text{in }\ k\hand{$[$}t\hand{$]$}/(\hand{$t^3$}), \end{array} \]
where a,b are parameters in the commutative ring k,  the group structure
being multiplication.  These parameters define in an evident way \over{\unsure{a}}{a} struc-%
ture of an extension of $\underline{G}_a$ by $\underline{G}_a$ upon $W_2$ ,  a splitting of which would
correspond to a group homomorphism
\[ \underline{G}_a \hand{$\longrightarrow$} W_2 \ , \quad a \hand{$\longmapsto$} 1 + at + P(a)t^2 \quad , \]
with
\[ P \hand{$\in$}\ Z\hand{$[$}t\hand{$]$} \ . \]
Expressing compatibility with the group laws gives the condition
\[ P(a+a') = P(a) + P(a') + aa' \quad , \]
which has\hand{,} as unique solutions  in  $\underline{Q}$\hand{$[$}t\hand{$]$}\hand{,} expressions
\[ P(t) = t(t-1)/2 \ + \ ct \]
with c in $\underline{Q}$ ,  none of which has coefficients in $\underline{Z}$ . \deleted{\unsure{(}}This argument shows
in fact that for given ring k, $(W_2)_k$ is a split extension iff 2 is inver-%
tible in k,  in which case a splitting is given by
\[ a \hand{$\longmapsto$} 1 + at +(a(a-1)/2)t \quad . \]

This example brings near one plausible ``reason'' why the expected
comparison statement about discrete and schematic linearization could not
reasonably hold true,  and in particular\inserted{why} we should'nt expect discrete and
schematic Hochs\over{\unsure{c}}{ch}ild cohomology (for group schemes over $\underline{Z}$  such as $\underline{G}_a$

\margin{59}
